\documentclass[12pt, a4paper]{article}
\usepackage[utf8]{inputenc}
\usepackage[T1]{fontenc}
\usepackage[french]{babel}
\usepackage[margin=2.5cm]{geometry}
\usepackage{amsmath, amssymb}
\usepackage{listings}
\usepackage{xcolor}
\usepackage{tcolorbox}
\tcbuselibrary{theorems, skins, breakable}
\usepackage{booktabs}
\usepackage{fancyhdr}
\usepackage{hyperref}
\usepackage{tikz}
\usetikzlibrary{arrows.meta, positioning, shapes.geometric}

\definecolor{suva_blue}{RGB}{30, 100, 180}
\definecolor{code_bg}{RGB}{245, 247, 250}
\definecolor{code_kw}{RGB}{0, 80, 200}
\definecolor{code_str}{RGB}{180, 40, 40}
\definecolor{code_cmt}{RGB}{80, 130, 80}
\definecolor{warn_bg}{RGB}{255, 248, 220}
\definecolor{success_bg}{RGB}{230, 255, 235}
\definecolor{sol_green}{RGB}{0, 120, 80}
\definecolor{sol_bg}{RGB}{240, 255, 245}

\lstdefinestyle{solidity}{
  language=C++,
  basicstyle=\ttfamily\small,
  keywordstyle=\color{code_kw}\bfseries,
  stringstyle=\color{code_str},
  commentstyle=\color{code_cmt}\itshape,
  backgroundcolor=\color{sol_bg},
  frame=single, rulecolor=\color{sol_green!40},
  numbers=left, numberstyle=\tiny\color{gray},
  breaklines=true, tabsize=4, showstringspaces=false,
  morekeywords={uint256, uint8, bytes, bytes32, bool, address, memory,
    storage, calldata, external, view, pure, returns, mapping, struct,
    event, emit, require, revert, modifier, payable, public, private,
    internal, virtual, override, assembly, let, mstore, mload},
}
\lstdefinestyle{python}{
  language=Python,
  basicstyle=\ttfamily\small,
  keywordstyle=\color{code_kw}\bfseries,
  stringstyle=\color{code_str},
  commentstyle=\color{code_cmt}\itshape,
  backgroundcolor=\color{code_bg},
  frame=single, rulecolor=\color{gray!40},
  numbers=left, numberstyle=\tiny\color{gray},
  breaklines=true, tabsize=4, showstringspaces=false,
}
\lstdefinestyle{bash}{
  basicstyle=\ttfamily\small\color{white},
  backgroundcolor=\color{gray!80},
  frame=single, rulecolor=\color{gray!80},
  breaklines=true,
}
\lstdefinestyle{output}{
  basicstyle=\ttfamily\small,
  backgroundcolor=\color{gray!10},
  frame=single, rulecolor=\color{gray!50},
  breaklines=true,
}

\newtcolorbox{defbox}[2][]{colback=success_bg,colframe=green!50!black,
  fonttitle=\bfseries,title={Définition : #2},#1}
\newtcolorbox{exercice}[2][]{colback=white,colframe=suva_blue!60,
  fonttitle=\bfseries,title={Exercice #2},#1,breakable}
\newtcolorbox{solution}[1][]{colback=gray!5,colframe=gray!40,
  fonttitle=\bfseries\itshape,title={Solution},#1,breakable}
\newtcolorbox{importantbox}[1][]{colback=suva_blue!10,colframe=suva_blue,
  fonttitle=\bfseries,title={Connexion avec le projet},#1}
\newtcolorbox{warningbox}[1][]{colback=warn_bg,colframe=orange!80!black,
  fonttitle=\bfseries,title={Attention},#1}
\newtcolorbox{gasbox}[1][]{colback=orange!5,colframe=orange!80,
  fonttitle=\bfseries,title={Gas cost},#1}

\pagestyle{fancy}
\fancyhf{}
\fancyhead[L]{\color{suva_blue}\textbf{SuVa --- Chapitre 6}}
\fancyhead[R]{\color{gray}Ethereum / Solidity / ERC-4337}
\fancyfoot[L]{\footnotesize\color{gray}Mohamed Lamine MBENGUE}
\fancyfoot[C]{\thepage}
\fancyfoot[R]{\footnotesize\color{gray}portfolio-alamine.web.app}
\renewcommand{\headrulewidth}{0.4pt}
\renewcommand{\footrulewidth}{0.3pt}

% ══════════════════════════════════════════════════════════════════════
\begin{document}

\begin{center}
  {\color{suva_blue}\rule{\linewidth}{2pt}}\\[0.5cm]
  {\Huge\bfseries\color{suva_blue} Chapitre 6}\\[0.3cm]
  {\large\bfseries Ethereum, Solidity, Gas et ERC-4337}\\[0.2cm]
  {\small\color{gray} La partie on-chain du projet SuVa}\\[0.3cm]
  {\color{suva_blue}\rule{\linewidth}{2pt}}
\end{center}

\vspace{0.3cm}
\begin{tcolorbox}[colback=suva_blue!5, colframe=suva_blue!40,
  left=6pt, right=6pt, top=4pt, bottom=4pt]
\small
\begin{tabular}{ll}
  \textbf{Auteur}      & Mohamed Lamine MBENGUE \\
  \textbf{Spécialités} & Sécurité · DevOps · Dev Full Stack Web \& Mobile \\
  \textbf{Contact}     & +221 78 178 44 36 · mbenguemohamedlamine2022@gmail.com \\
  \textbf{Portfolio}   & \href{https://portfolio-alamine.web.app}{\texttt{portfolio-alamine.web.app}} \\
\end{tabular}
\end{tcolorbox}
\vspace{0.3cm}

\tableofcontents
\newpage

%══════════════════════════════════════════════════════════════════════
\section{L'EVM --- Ethereum Virtual Machine}
%══════════════════════════════════════════════════════════════════════

\subsection{Fonctionnement de base}

L'EVM est une machine à pile (\emph{stack machine}) qui exécute le bytecode
des smart contracts. Chaque opération (opcode) a un coût en \textbf{gas}.

\begin{defbox}{Gas}
Le \textbf{gas} est l'unité de mesure du coût de calcul sur Ethereum.
\begin{itemize}
  \item Chaque opcode coûte un certain nombre de gas
  \item Le gas est payé en ETH (gas\_used $\times$ gas\_price)
  \item La limite de gas par bloc est $\sim 30\,000\,000$ gas
\end{itemize}
\end{defbox}

\subsection{Opcodes importants pour ce projet}

\begin{center}
\begin{tabular}{llr}
\toprule
\textbf{Opcode} & \textbf{Opération} & \textbf{Gas} \\
\midrule
\texttt{ADD} & Addition 256-bit & 3 \\
\texttt{MUL} & Multiplication 256-bit & 5 \\
\texttt{DIV} & Division 256-bit & 5 \\
\texttt{ADDMOD} & (a+b) mod n & 8 \\
\texttt{MULMOD} & (a*b) mod n & 8 \\
\texttt{KECCAK256} & Hash Keccak-256 & 30 + 6/word \\
\texttt{MLOAD} & Lecture mémoire & 3 \\
\texttt{MSTORE} & Écriture mémoire & 3 \\
\texttt{EXTCODECOPY} & Copier bytecode & 100 + copie \\
\bottomrule
\end{tabular}
\end{center}

\begin{gasbox}
Comparaison des coûts de vérification post-quantique~:
\begin{center}
\begin{tabular}{lr}
\toprule
\textbf{Schéma} & \textbf{Gas} \\
\midrule
ECDSA (actuel Ethereum) & $\sim$ 3\,000 \\
\textbf{Falcon (NIST FIPS 206)} & $\sim$ \textbf{350\,000} \\
ETHDilithium (Keccak) & $\sim$ 6\,600\,000 \\
Dilithium standard (SHAKE) & $\sim$ 13\,500\,000 \\
SPHINCS+ & $\sim$ 50--500\,000\,000 \\
Limite par bloc & 30\,000\,000 \\
\bottomrule
\end{tabular}
\end{center}
Falcon est $\sim 20\times$ moins cher que Dilithium et tient bien dans la limite de bloc.
\end{gasbox}

%══════════════════════════════════════════════════════════════════════
\section{Solidity --- Bases essentielles}
%══════════════════════════════════════════════════════════════════════

\subsection{Types de données utiles dans ce projet}

\begin{lstlisting}[style=solidity]
// Types de base
uint256 a = 42;              // entier 256-bit non signe
bytes32 h;                   // 32 bytes fixes (hash)
bytes memory data;           // tableau de bytes dynamique
bool ok = true;

// Tableaux
uint256[] memory arr;                 // tableau dynamique
uint256[256] memory poly;             // polynome (256 coefficients)
uint256[][] memory vector;            // vecteur de polynomes
uint256[][][] memory matrix;          // matrice de polynomes

// Structures (utilises pour PubKey et Signature)
struct PubKey {
    bytes32 rho;            // graine de la matrice A
    bytes tr;               // hash de la cle publique
    uint256[][] t1;         // vecteur t1 (k polynomes)
}

struct Signature {
    bytes32[] c_tilde;      // hash du challenge (32 bytes)
    uint256[][] z;          // vecteur z (l polynomes)
    uint256[][] h;          // hint (k polynomes de bits)
}
\end{lstlisting}

\subsection{Les zones mémoire}

\begin{defbox}{Memory vs Storage vs Calldata}
\begin{itemize}
  \item \textbf{storage}~: donnees persistantes sur la blockchain (très cher)
  \item \textbf{memory}~: donnees temporaires pendant l'execution (moderé)
  \item \textbf{calldata}~: arguments en lecture seule d'un appel (le moins cher)
\end{itemize}
\end{defbox}

\begin{importantbox}
Dans les contrats Dilithium, presque tout est en \texttt{memory}
car la vérification est stateless (pas de stockage permanent).
Chaque appel à \texttt{verify()} recommence de zéro.
\end{importantbox}

\subsection{Fonctions et visibilité}

\begin{lstlisting}[style=solidity]
// Fonction de verification (interface externe)
function verify(
    PubKey memory pk,
    bytes memory message,
    Signature memory signature
) external view returns (bool) {
    // external : appelee depuis l'exterieur
    // view     : ne modifie pas l'etat
    // returns  : type de retour
    ...
}

// Fonction interne (helper)
function compute_c_ntt(
    Signature memory sig
) internal pure returns (uint256[] memory) {
    // internal : uniquement depuis ce contrat
    // pure     : pas de lecture d'etat non plus
    ...
}
\end{lstlisting}

%══════════════════════════════════════════════════════════════════════
\section{Yul --- EVM Assembly dans Solidity}
%══════════════════════════════════════════════════════════════════════

\subsection{Pourquoi Yul dans ce projet ?}

Le NTT (\texttt{ZKNOX\_NTT.sol}) est écrit entièrement en Yul pour
maximiser l'efficacité du gas. En Solidity pur, le compilateur ajoute
des vérifications et des copies inutiles. Yul donne un contrôle total.

\subsection{Syntaxe de base Yul}

\begin{lstlisting}[style=solidity]
// Dans Solidity, on utilise assembly { ... } pour du Yul inline
function mulmod_example(uint256 a, uint256 b, uint256 m)
    internal pure returns (uint256 result) {
    assembly {
        // Variables locales
        let x := mulmod(a, b, m)   // (a * b) mod m

        // Operations sur la memoire
        let ptr := mload(0x40)     // Lire le free memory pointer
        mstore(ptr, x)             // Ecrire x a l'adresse ptr

        // Affectation du resultat
        result := x
    }
}

// NTT butterfly step (simplifie)
function butterfly(uint256 u, uint256 v, uint256 w, uint256 q)
    internal pure returns (uint256 u_out, uint256 v_out) {
    assembly {
        // t = w * v mod q
        let t := mulmod(w, v, q)
        // u_out = (u + t) mod q
        u_out := addmod(u, t, q)
        // v_out = (u - t) mod q
        v_out := addmod(u, sub(q, t), q)
    }
}
\end{lstlisting}

\subsection{extcodecopy --- lire les twiddle factors}

\begin{lstlisting}[style=solidity]
// Dans ZKNOX_NTT.sol
// Les twiddle factors sont stockes comme bytecode d'un contrat
// (deploye par ZKNOX_dilithium_deploy.sol)

function load_twiddle_factors(address twiddle_contract)
    internal view returns (uint256[] memory) {
    assembly {
        // Taille du bytecode (= nb_twiddles * 32 bytes)
        let size := extcodesize(twiddle_contract)

        // Allouer la memoire
        let ptr := mload(0x40)
        mstore(0x40, add(ptr, size))

        // Copier le bytecode (qui contient les constantes)
        extcodecopy(twiddle_contract, ptr, 0, size)
    }
}
// Avantage : extcodecopy est moins cher que lire depuis storage
// Inconvenient : les constantes doivent etre deployees d'abord
\end{lstlisting}

%══════════════════════════════════════════════════════════════════════
\section{Architecture des contrats Solidity dans le projet}
%══════════════════════════════════════════════════════════════════════

\begin{center}
\begin{tikzpicture}[
  block/.style={draw=sol_green, rounded corners, fill=sol_bg,
    minimum width=4cm, minimum height=0.8cm, text width=3.8cm,
    align=center, font=\small},
  arrow/.style={->, thick, sol_green!80}
]
  \node[block] (entry) {\texttt{ZKNOX\_dilithium.sol}\\{\tiny entry point, verify()}};
  \node[block, below left=1cm and 0.5cm of entry]
    (core) {\texttt{ZKNOX\_dilithium\_core.sol}\\{\tiny core\_1(), core\_2()}};
  \node[block, below right=1cm and 0.5cm of entry]
    (eth) {\texttt{ZKNOX\_ethdilithium.sol}\\{\tiny variant Keccak}};
  \node[block, below=1cm of core]
    (ntt) {\texttt{ZKNOX\_NTT.sol}\\{\tiny NTT en Yul}};
  \node[block, below left=1cm and 0.5cm of ntt]
    (nttdil) {\texttt{ZKNOX\_NTT\_dilithium.sol}\\{\tiny NTTFW, NTTINV}};
  \node[block, below right=1cm and 0.5cm of ntt]
    (utils) {\texttt{ZKNOX\_dilithium\_utils.sol}\\{\tiny constantes, structs}};
  \node[block, below=1cm of core, xshift=3.5cm]
    (hint) {\texttt{ZKNOX\_hint.sol}\\{\tiny useHint()}};
  \node[block, right=1cm of entry]
    (shake) {\texttt{ZKNOX\_shake.sol}\\{\tiny SHAKE-256 Solidity}};
  \node[block, below=0.5cm of shake]
    (keccak) {\texttt{ZKNOX\_keccak\_prng.sol}\\{\tiny Keccak PRNG}};

  \draw[arrow] (entry) -- (core);
  \draw[arrow] (entry) -- (eth);
  \draw[arrow] (core) -- (ntt);
  \draw[arrow] (ntt) -- (nttdil);
  \draw[arrow] (ntt) -- (utils);
  \draw[arrow] (core) -- (hint);
  \draw[arrow] (entry) -- (shake);
  \draw[arrow] (eth) -- (keccak);
\end{tikzpicture}
\end{center}

\subsection{ZKNOX\_dilithium.sol --- le contrat principal}

\begin{lstlisting}[style=solidity]
// ZKNOX_dilithium.sol (simplifie)
pragma solidity ^0.8.25;

import "./ZKNOX_dilithium_core.sol";
import "./ZKNOX_shake.sol";
import "./ZKNOX_NTT_dilithium.sol";

contract ZKNOX_dilithium is ZKNOX_dilithium_core {
    // Adresses des contrats de twiddle factors (deployes)
    address immutable apsirev;    // NTT forward
    address immutable apsiInvrev; // NTT inverse

    constructor(address _apsirev, address _apsiInvrev) {
        apsirev = _apsirev;
        apsiInvrev = _apsiInvrev;
    }

    function verify(
        PubKey memory pk,
        bytes memory msgs,
        Signature memory signature
    ) external view returns (bool) {
        // STEP 1: Valider la structure de la signature
        (uint256 norm_h, uint256[][] memory h, uint256[][] memory z) =
            dilithium_core_1(signature);

        if (norm_h > OMEGA) return false;

        // Verifier que z est dans [-gamma_1+beta, gamma_1-beta]
        for (uint i = 0; i < L; i++) {
            for (uint j = 0; j < N; j++) {
                if (z[i][j] > GAMMA_1_MINUS_BETA &&
                    Q - z[i][j] > GAMMA_1_MINUS_BETA) {
                    return false;
                }
            }
        }

        // STEP 2: Reconstruire c en domaine NTT
        uint256[] memory c_ntt = compute_c_ntt(signature);

        // STEP 3: t1 * 2^d en NTT
        uint256[][] memory t1_new = ZKNOX_Expand_Vec(pk.t1);
        for (uint i = 0; i < K; i++) {
            for (uint j = 0; j < N; j++) {
                t1_new[i][j] <<= D;  // shift par 2^d
            }
            t1_new[i] = ZKNOX_NTTFW(t1_new[i], apsirev);
        }

        // STEP 4: w' = UseHint(h, A*z - c*t1)
        bytes memory w_prime = dilithium_core_2(
            apsirev, apsiInvrev, pk, z, c_ntt, h, t1_new
        );

        // STEP 5: Verification du hash final
        ctx_shake memory ctx;
        ctx = shake_update(ctx, pk.tr);
        ctx = shake_update(ctx, msgs);
        bytes memory mu = shake_digest(ctx, 64);

        ctx_shake memory ctx2;
        ctx2 = shake_update(ctx2, mu);
        ctx2 = shake_update(ctx2, w_prime);
        bytes32 final_hash = bytes32(shake_digest(ctx2, 32));

        for (uint i = 0; i < 32; i++) {
            if (signature.c_tilde[i] != uint8(final_hash[i]))
                return false;
        }
        return true;
    }
}
\end{lstlisting}

%══════════════════════════════════════════════════════════════════════
\section{ERC-4337 --- Account Abstraction}
%══════════════════════════════════════════════════════════════════════

\subsection{Pourquoi ERC-4337 ?}

Ethereum valide nativement les transactions via ECDSA.
Pour utiliser une signature post-quantique, on ne peut pas changer
le protocole Ethereum. ERC-4337 permet de contourner ça via des
\textbf{smart contract wallets}.

\subsection{Composants ERC-4337}

\begin{center}
\begin{tikzpicture}[
  block/.style={draw=suva_blue, rounded corners, fill=blue!5,
    minimum width=3.5cm, minimum height=1cm, text width=3.3cm,
    align=center, font=\small},
  arrow/.style={->, thick, suva_blue}
]
  \node[block] (user) {Utilisateur\\{\tiny Génère UserOperation}};
  \node[block, right=1.5cm of user] (mempool)
    {UserOperation\\Mempool\\{\tiny File d'attente}};
  \node[block, right=1.5cm of mempool]
    (bundler) {Bundler\\{\tiny Agrège les ops}};
  \node[block, below=1.5cm of bundler]
    (entry) {EntryPoint\\Contract\\{\tiny Contrat central}};
  \node[block, below=1.5cm of entry]
    (wallet) {Smart Wallet\\(SuVa)\\{\tiny validateUserOp()}};
  \node[block, right=1.5cm of entry]
    (paymaster) {Paymaster\\{\tiny Paie le gas}};

  \draw[arrow] (user) -- (mempool);
  \draw[arrow] (mempool) -- (bundler);
  \draw[arrow] (bundler) -- (entry);
  \draw[arrow] (entry) -- (wallet);
  \draw[arrow] (entry) -- (paymaster);
\end{tikzpicture}
\end{center}

\subsection{Le Smart Wallet SuVa}

\begin{lstlisting}[style=solidity]
// Contrat SuVa Wallet (a construire !)
// Ce contrat sera implemente dans le futur du projet

pragma solidity ^0.8.25;

import "@account-abstraction/contracts/interfaces/IAccount.sol";
import "./ZKNOX_ethdilithium.sol";

contract SuVaWallet is IAccount {
    // Cle publique post-quantique de l'utilisateur
    PubKey public quantum_pk;

    // Verifieur ETHDilithium
    ZKNOX_ethdilithium immutable verifier;

    constructor(
        PubKey memory _quantum_pk,
        address _verifier
    ) {
        quantum_pk = _quantum_pk;
        verifier = ZKNOX_ethdilithium(_verifier);
    }

    // Fonction appelee par EntryPoint pour valider une transaction
    function validateUserOp(
        UserOperation calldata userOp,
        bytes32 userOpHash,
        uint256 missingAccountFunds
    ) external returns (uint256 validationData) {
        // La signature dans userOp contient la signature post-quantique
        Signature memory sig = abi.decode(
            userOp.signature,
            (Signature)
        );

        // Verifier avec ETHDilithium
        bool valid = verifier.verify(
            quantum_pk,
            abi.encodePacked(userOpHash),
            sig
        );

        // 0 = valide, 1 = invalide
        return valid ? 0 : 1;
    }
}
\end{lstlisting}

\begin{importantbox}
Le flow complet pour une transaction SuVa~:
\begin{enumerate}
  \item L'utilisateur génère une signature Dilithium/Falcon off-chain (Python/SDK)
  \item La signature est encodée dans \texttt{UserOperation.signature}
  \item Le Bundler soumet la \texttt{UserOperation} à l'EntryPoint
  \item L'EntryPoint appelle \texttt{validateUserOp()} du Smart Wallet
  \item Le Smart Wallet appelle \texttt{verify()} du contrat ZKNOX
  \item Si la vérification passe, la transaction est exécutée
\end{enumerate}
\end{importantbox}

%══════════════════════════════════════════════════════════════════════
\section{Foundry --- Tester les contrats}
%══════════════════════════════════════════════════════════════════════

\subsection{Structure d'un test Foundry}

\begin{lstlisting}[style=solidity]
// test/ZKNOX_dilithium.t.sol (genere par Python)
pragma solidity ^0.8.25;

import "forge-std/Test.sol";
import "../src/ZKNOX_dilithium.sol";

contract TestDilithium is Test {
    ZKNOX_dilithium verifier;

    // Donnees de test (generees par generate_dilithium_test_vectors.py)
    bytes constant PUBLIC_KEY = hex"abc123..."; // 1312 bytes
    bytes constant MESSAGE = "We are ZKNox.";
    bytes constant SIGNATURE = hex"def456..."; // 2420 bytes

    function setUp() public {
        // Deployer le contrat (et les twiddle factors)
        // ...
        verifier = new ZKNOX_dilithium(...);
    }

    function test_verify_valid_signature() public {
        // Decoder les donnees
        PubKey memory pk = abi.decode(PUBLIC_KEY, (PubKey));
        Signature memory sig = abi.decode(SIGNATURE, (Signature));

        // Verifier
        bool result = verifier.verify(pk, MESSAGE, sig);
        assertTrue(result, "La signature valide doit passer");
    }

    function test_verify_wrong_message() public {
        PubKey memory pk = abi.decode(PUBLIC_KEY, (PubKey));
        Signature memory sig = abi.decode(SIGNATURE, (Signature));

        bool result = verifier.verify(pk, "Wrong message", sig);
        assertFalse(result, "Mauvais message doit echouer");
    }

    function test_gas_cost() public {
        // Foundry mesure automatiquement le gas si on utilise
        // vm.startSnapshotGas / vm.stopSnapshotGas
        uint256 gasBefore = gasleft();
        // ... (verification)
        uint256 gasUsed = gasBefore - gasleft();
        emit log_uint(gasUsed);  // ~6.6M pour ETHDilithium
    }
}
\end{lstlisting}

\subsection{Commandes Foundry essentielles}

\begin{lstlisting}[style=bash]
# Compiler les contrats
forge build

# Lancer tous les tests
forge test -vv

# Tester seulement un contrat specifique
forge test --match-contract TestDilithium -vv

# Tester une seule fonction
forge test --match-test test_verify_valid_signature -vvv

# Rapport gas detaille
forge test --gas-report

# Profil lite (plus rapide)
FOUNDRY_PROFILE=lite forge test -vv
\end{lstlisting}

%══════════════════════════════════════════════════════════════════════
\section{Exercices}
%══════════════════════════════════════════════════════════════════════

\begin{exercice}{1 --- Lecture de contrat Solidity}
Ouvre \texttt{src/ZKNOX\_dilithium\_utils.sol}.
\begin{enumerate}[label=\alph*)]
  \item Identifie toutes les constantes Dilithium définies.
        Vérifie qu'elles correspondent aux valeurs du Chapitre 5.
  \item Quelle est la structure \texttt{PubKey} ? Quels champs contient-elle ?
  \item Quelle est la structure \texttt{Signature} ? Quels champs contient-elle ?
  \item Trouve la fonction \texttt{VECMULMOD} --- que fait-elle ?
\end{enumerate}
\end{exercice}

\begin{exercice}{2 --- Coûts gas}
\begin{enumerate}[label=\alph*)]
  \item Calcule le coût en ETH d'une vérification ETHDilithium (6.6M gas)
        si le gas price est 20 Gwei et 1 ETH = 3000\$.
  \item Même calcul pour Falcon (350K gas).
  \item Sur un L2 comme Arbitrum (gas 10x moins cher), quel serait le coût ?
\end{enumerate}
\end{exercice}

\begin{solution}
\begin{lstlisting}[style=python]
# Parametres
gas_ethdilithium = 6_600_000
gas_falcon = 350_000
gas_price_gwei = 20
eth_usd = 3000

# ETHDilithium
cost_eth = gas_ethdilithium * gas_price_gwei * 1e-9  # ETH
cost_usd = cost_eth * eth_usd
print(f"ETHDilithium : {cost_eth:.4f} ETH = ${cost_usd:.2f}")
# ~ 0.132 ETH = $396

# Falcon
cost_eth_f = gas_falcon * gas_price_gwei * 1e-9
cost_usd_f = cost_eth_f * eth_usd
print(f"Falcon : {cost_eth_f:.4f} ETH = ${cost_usd_f:.2f}")
# ~ 0.007 ETH = $21

# L2 (10x moins cher)
print(f"Falcon L2 : ${cost_usd_f / 10:.2f}")
# ~ $2.10
\end{lstlisting}
\end{solution}

\begin{exercice}{3 --- ERC-4337}
\begin{enumerate}[label=\alph*)]
  \item Quel est le rôle du \texttt{Bundler} dans ERC-4337 ?
  \item Pourquoi le Smart Wallet paie-t-il le gas (via le Paymaster)
        plutôt que l'utilisateur directement ?
  \item Pour SuVa, quel est l'avantage de ERC-4337 vs modifier le
        protocole Ethereum directement ?
\end{enumerate}
\end{exercice}

\begin{exercice}{4 --- Installer et compiler le projet}
\begin{enumerate}[label=\alph*)]
  \item Installe Foundry et compile les contrats : \texttt{forge build}.
  \item Y a-t-il des erreurs de compilation ? Si oui, lesquelles ?
  \item Lance \texttt{FOUNDRY\_PROFILE=lite forge test -vv} et observe la sortie.
\end{enumerate}
\end{exercice}

\end{document}
